\chapter{Conclusion}
\label{chap:Conclusion}

\section{Summary}
\label{sec:Summary}

This thesis introduced provider-selective distillation, a curriculum over providers and
domains in which a small student chooses whom to learn from, about what, and when. A
preliminary study of its first version, with priced black-box providers, found that
token-level supervision helps only when the tokenizer is shared, that price rather than
provider strength decides value under a budget, and that a selector which learns where
progress happens matches the best single provider without being told which one it is.
Learning was confined to Countdown and was unstable. These findings led to the refined
problem: no budget, white-box providers of the student's family, and learning curves as
the measure. For it, the thesis proposes a training objective that combines a
ground-truth reward with the selected provider's token probabilities, and a selector
that scores each arm by its learned gain per round. The evaluation of
Chapter~\ref{chap:Plan} will test whether the objective is stable and makes progress on
every domain, whether an adaptive curriculum is faster than a fixed one, and whether a
student trained on several providers passes its teachers.

\section{Limitations}
\label{sec:Limitations}

Every result in this report comes from the budgeted, black-box version of the problem;
the refined problem has not yet been run. In that study the providers answer from a
fixed cache of four responses per prompt, prices are a compute-based proxy rather than
real API prices, and the student learns almost only on Countdown, so most of the
evidence comes from one domain. Three seeds give wide intervals. The student's size was
not varied, and tasks are single-turn with checkable answers. The limitations of the
design for the refined problem are listed in Section~\ref{imp:sec:limitations}.

\section{Future Work}
\label{sec:Future Work}

Beyond the planned evaluation, three directions follow from this work. The first is to
bring back providers of other model families through a text-only signal, alongside the
white-box ones, so that the curriculum also chooses the form of supervision. The second
is a selector that needs no labelled selection set, using a signal available during
training such as the size of the advantages. The third is to extend the tasks from
single-turn problems with checkable answers to multi-turn agentic tasks.
